% =========================================================================
% UNIT 7: INFERENCE FOR QUANTITATIVE DATA - MEANS
% =========================================================================
\chapter{Unit 7: Inference for Quantitative Data: Means}

\section{Unit Overview \& Exam Weight}
Unit 7 accounts for \textbf{10\% to 18\%} of the AP Statistics Exam. Whenever dealing with quantitative variables where the true population standard deviation $\sigma$ is unknown (which is practically 100\% of real-world scenarios), we use the \textbf{$t$-distribution} discovered by William Sealy Gosset (writing under the pseudonym "Student"). You will master 1-sample $t$-procedures, paired $t$-procedures, and 2-sample $t$-procedures.

\subsection*{Key Unit Vocabulary}
\begin{description}[leftmargin=1.5em, style=nextline]
  \item[Student's $t$-Distribution:] A family of bell-shaped, symmetric density curves defined by degrees of freedom ($df$). It has heavier tails and more probability in the extremes than the standard normal distribution to account for the extra uncertainty of estimating $\sigma$ with $s_x$. As $df \to \infty$, the $t$-distribution approaches $N(0, 1)$.
  \item[Degrees of Freedom ($df$):] For a 1-sample $t$-distribution, $df = n - 1$.
  \item[Standard Error of the Mean:] $\text{SE}(\bar{x}) = \frac{s_x}{\sqrt{n}}$.
  \item[1-Sample $t$-Interval:] An interval estimating an unknown population mean $\mu$: $\bar{x} \pm t^* \frac{s_x}{\sqrt{n}}$.
  \item[1-Sample $t$-Test:] A hypothesis test evaluating claims about a population mean $\mu$: $t = \frac{\bar{x} - \mu_0}{s_x / \sqrt{n}}$.
  \item[Paired Data (Matched Pairs):] Quantitative data resulting from paired observations on the same individuals (e.g., pre-test vs. post-test) or pairs of similar individuals. Analyzed by taking differences ($d = x_1 - x_2$) and running 1-sample $t$-procedures on the differences ($\mu_d, \bar{x}_d, s_d$).
  \item[Two-Sample $t$-Procedures:] Inference comparing two independent population means $\mu_1 - \mu_2$. Never pool variances for means!
  \item[Robustness:] An inference procedure is robust if the stated confidence level or $p$-value remains approximately accurate even when a condition (such as normality) is moderately violated. The $t$-procedures are robust against non-normality when $n \ge 30$.
\end{description}

\section{Core Concept Lessons}

\begin{lessonbox}[Lesson 7.1: The $t$-Distribution vs. Standard Normal ($z$)]
Why can we NOT use $z$ when working with sample means?
Because the true population standard deviation $\sigma$ is almost never known! When we replace $\sigma$ with the sample standard deviation $s_x$, we introduce additional sampling variability. To compensate for this extra variation, the $t$-distribution spreads out wider, with fatter tails.
\begin{itemize}[leftmargin=1.2em]
  \item Bell-shaped and symmetric about 0.
  \item Area under curve = 1.
  \item Controlled by degrees of freedom: $df = n - 1$.
  \item As $df$ increases, $s_x$ becomes a more precise estimate of $\sigma$, and the $t$-distribution converges to the standard normal $Z$ curve!
\end{itemize}
\end{lessonbox}

\begin{lessonbox}[Lesson 7.2: One-Sample $t$-Interval for a Population Mean $\mu$]
\textbf{Confidence Interval Formula:}
\begin{equation}
  \bar{x} \pm t^* \left(\frac{s_x}{\sqrt{n}}\right)
\end{equation}
where $t^*$ is the critical value with $df = n - 1$.
\textbf{Mandatory Conditions Check:}
\begin{enumerate}[leftmargin=1.5em]
  \item \textbf{Random:} Data comes from an SRS or randomized experiment.
  \item \textbf{10\% Condition:} $n \le 0.10N$ when sampling without replacement.
  \item \textbf{Normal / Large Sample Condition:}
  \begin{itemize}
    \item If the population is stated to be \textbf{Normal}, the condition is met.
    \item If $n \ge 30$, the Central Limit Theorem guarantees the sampling distribution is approximately normal.
    \item If $n < 30$, you MUST graph the sample data (dotplot, stemplot, or boxplot) and confirm there are \textbf{no strong skewness or outliers}!
  \end{itemize}
\end{enumerate}
\end{lessonbox}

\begin{lessonbox}[Lesson 7.3: Paired $t$-Procedures (Matched Pairs Design)]
\textbf{Critical Distinction:} When two sets of quantitative measurements are collected on the \textbf{same individuals} (e.g., before-treatment and after-treatment) or matched pairs, they are NOT independent!
\begin{itemize}[leftmargin=1.2em]
  \item Compute individual differences for each pair: $d_i = x_{1i} - x_{2i}$.
  \item Treat the list of differences as a \textbf{single sample} of size $n$ with sample mean difference $\bar{x}_d$ and sample standard deviation of differences $s_d$.
  \item Hypotheses: $H_0: \mu_d = 0$ versus $H_a: \mu_d > 0, \mu_d < 0,$ or $\mu_d \neq 0$.
  \item Test Statistic:
  \begin{equation}
    t = \frac{\bar{x}_d - 0}{s_d / \sqrt{n}}, \qquad df = n - 1
  \end{equation}
  \item Confidence Interval: $\bar{x}_d \pm t^* \frac{s_d}{\sqrt{n}}$.
\end{itemize}
\end{lessonbox}

\begin{lessonbox}[Lesson 7.4: Two-Sample $t$-Procedures ($\mu_1 - \mu_2$)]
Comparing the means of two independent groups:
\begin{itemize}[leftmargin=1.2em]
  \item \textbf{Confidence Interval:}
  \begin{equation}
    (\bar{x}_1 - \bar{x}_2) \pm t^* \sqrt{\frac{s_1^2}{n_1} + \frac{s_2^2}{n_2}}
  \end{equation}
  \item \textbf{Test Statistic ($H_0: \mu_1 - \mu_2 = 0$):}
  \begin{equation}
    t = \frac{(\bar{x}_1 - \bar{x}_2) - 0}{\sqrt{\frac{s_1^2}{n_1} + \frac{s_2^2}{n_2}}}
  \end{equation}
  \item \textbf{Degrees of Freedom:} Use technology (Welch-Satterthwaite approximation) or the conservative formula: $df = \min(n_1 - 1, n_2 - 1)$.
  \item \textbf{Important Rule:} \textbf{NEVER POOL} variances for two-sample $t$-procedures on the AP exam! Always select \texttt{Pooled: NO} on your TI-84.
\end{itemize}
\end{lessonbox}

\section{Worked Step-by-Step Examples}

\subsection*{Example 7.1: One-Sample $t$-Test for a Mean (PHANTOM)}
\textbf{Problem:} A consumer agency tests whether a brand of AAA batteries lasts less than the advertised mean lifetime of 30 hours. A random sample of $n = 16$ batteries has a mean lifetime of $\bar{x} = 28.5\text{ hours}$ with a standard deviation of $s = 2.4\text{ hours}$. A dotplot of the 16 values shows no outliers or strong skewness. Conduct a test at $\alpha = 0.05$.

\textbf{Solution:}
\begin{itemize}[leftmargin=1.2em]
  \item \textbf{P:} Let $\mu = \text{the true mean lifetime of all AAA batteries of this brand (in hours)}$.
  \item \textbf{H:} $H_0: \mu = 30$ versus $H_a: \mu < 30$.
  \item \textbf{A:}
  1. \textit{Random:} Stated random sample of 16 batteries.
  2. \textit{10\% Condition:} $n = 16 \le 0.10(\text{all batteries produced})$.
  3. \textit{Normal/Large Sample:} Although $n = 16 < 30$, the problem states the dotplot of sample data shows no outliers or strong skewness. Therefore, $t$-procedures are safe to use.
  \item \textbf{N:} One-sample $t$-test for a population mean $\mu$.
  \item \textbf{T:} $df = n - 1 = 16 - 1 = 15$.
  \[ t = \frac{\bar{x} - \mu_0}{s / \sqrt{n}} = \frac{28.5 - 30}{2.4 / \sqrt{16}} = \frac{-1.5}{0.60} = -2.50 \]
  \item \textbf{O:} $p\text{-value} = P(t_{15} \le -2.50) = \mathbf{0.0123}$.
  \item \textbf{M:} Because the $p$-value ($0.0123$) is less than $\alpha = 0.05$, we reject $H_0$. There is convincing statistical evidence that the true mean lifetime of this battery brand is strictly less than 30 hours.
\end{itemize}

\section{TI-84 Plus CE Calculator Playbook}

\begin{calculatorbox}[TI-84 Playbook: Inference for Means]
\begin{itemize}[leftmargin=1.2em]
  \item \textbf{1-Sample $t$-Test:} Press \texttt{[STAT] $\rightarrow$ TESTS $\rightarrow$ 2:T-Test}.
  Select \texttt{Inpt: Stats} (or \texttt{Data} if raw list in \texttt{L1}). Enter $\mu_0, \bar{x}, s_x, n$, and direction of $H_a$.
  \item \textbf{1-Sample $t$-Interval:} Press \texttt{[STAT] $\rightarrow$ TESTS $\rightarrow$ 8:TInterval}. Enter stats, C-Level.
  \item \textbf{Paired $t$-Test:} Enter before-data in \texttt{L1}, after-data in \texttt{L2}. In \texttt{L3}, define \texttt{L1 - L2}. Run \texttt{2:T-Test} using \texttt{Data: L3} with $\mu_0 = 0$.
  \item \textbf{2-Sample $t$-Test:} Press \texttt{[STAT] $\rightarrow$ TESTS $\rightarrow$ 4:2-SampTTest}. Always set \texttt{Pooled: NO}!
  \item \textbf{2-Sample $t$-Interval:} Press \texttt{[STAT] $\rightarrow$ TESTS $\rightarrow$ 0:2-SampTInt}. Always set \texttt{Pooled: NO}!
\end{itemize}
\end{calculatorbox}

\section{Comprehensive Practice Problem Set}

\subsection*{Part I: 20 Multiple-Choice Questions}

\begin{enumerate}[leftmargin=1.5em]
  \item Why do we use the $t$-distribution instead of the standard normal distribution when constructing a confidence interval for a population mean?
  \begin{enumerate}[label=(\Alph*)]
    \item Because the sample size is always small.
    \item Because the population standard deviation $\sigma$ is unknown and estimated by $s$.
    \item Because the population is always non-normal.
    \item Because the mean is non-resistant.
    \item To avoid calculating degrees of freedom.
  \end{enumerate}
  \textbf{Answer: (B).} Estimating $\sigma$ with $s$ introduces additional variability, requiring $t$.

  \item As degrees of freedom increase, the $t$-distribution:
  \begin{enumerate}[label=(\Alph*)]
    \item Becomes more skewed. \item Approaches the standard normal distribution $Z$. \item Becomes wider with fatter tails. \item Has a mean greater than 0. \item Has a standard deviation of 0.
  \end{enumerate}
  \textbf{Answer: (B).} As $df \to \infty$, $t \to Z$.

  \item A random sample of $n = 25$ students has $\bar{x} = 75, s = 10$. For a 95\% confidence interval for $\mu$, what is $df$ and $t^*$?
  \begin{enumerate}[label=(\Alph*)]
    \item $df = 25, t^* = 1.960$ \item $df = 24, t^* = 2.064$ \item $df = 24, t^* = 1.960$ \item $df = 25, t^* = 2.060$ \item $df = 23, t^* = 2.500$
  \end{enumerate}
  \textbf{Answer: (B).} $df = n - 1 = 24$. From the $t$-table, $t^* = 2.064$.

  \item A matched pairs study measures cholesterol before and after a diet program for 20 patients. What degrees of freedom should be used for the hypothesis test?
  \begin{enumerate}[label=(\Alph*)]
    \item 38 \item 39 \item 19 \item 20 \item 18
  \end{enumerate}
  \textbf{Answer: (C).} Paired data is analyzed as a single sample of $n = 20$ differences $\implies df = 20 - 1 = 19$.

  \item If $n = 10$ and the sample data displays a severe outlier, is it appropriate to conduct a 1-sample $t$-test?
  \begin{enumerate}[label=(\Alph*)]
    \item Yes, because $t$-procedures are always robust.
    \item No, because for $n < 30$, the presence of an outlier violates the normality condition.
    \item Yes, if $\alpha = 0.01$.
    \item Yes, by the Central Limit Theorem.
    \item No, because $10 < 30$ is never allowed.
  \end{enumerate}
  \textbf{Answer: (B).} For small samples ($n < 30$), outliers invalidate $t$-procedures.

  \item A 95\% confidence interval for $\mu_1 - \mu_2$ is $(-4.5, -1.2)$. What can be concluded?
  \begin{enumerate}[label=(\Alph*)]
    \item $\mu_1 = \mu_2$.
    \item There is convincing evidence that $\mu_1 < \mu_2$ because the entire interval is negative (strictly below 0).
    \item There is no significant difference between the two means.
    \item $\mu_1 > \mu_2$.
    \item The test was two-sided.
  \end{enumerate}
  \textbf{Answer: (B).} All plausible values for $\mu_1 - \mu_2$ are negative $\implies \mu_1 < \mu_2$.

  \item In a two-sample $t$-test, why do AP Statistics standards mandate "Never Pool"?
  \begin{enumerate}[label=(\Alph*)]
    \item Pooling requires assuming the two population variances are equal ($\sigma_1^2 = \sigma_2^2$), an assumption rarely justified in practice.
    \item Pooling decreases power.
    \item Calculators cannot compute pooled tests.
    \item Pooling is allowed only for proportions.
    \item Pooling violates the 10\% condition.
  \end{enumerate}
  \textbf{Answer: (A).} Unequal variance $t$-tests (Welch) are far more robust and realistic.

  \item If a 90\% confidence interval for $\mu$ is $(45.2, 54.8)$, what was the sample mean $\bar{x}$?
  \begin{enumerate}[label=(\Alph*)]
    \item 45.2 \item 50.0 \item 54.8 \item 4.8 \item 9.6
  \end{enumerate}
  \textbf{Answer: (B).} Center of interval: $(45.2 + 54.8)/2 = 100/2 = 50.0$.

  \item A test of $H_0: \mu = 100$ vs. $H_a: \mu \neq 100$ produces $t = 2.45$ with $df = 15$. The $p$-value is:
  \begin{enumerate}[label=(\Alph*)]
    \item $P(t_{15} > 2.45)$ \item $2 \times P(t_{15} > 2.45)$ \item $P(t_{15} < 2.45)$ \item 0.05 \item 0.0245
  \end{enumerate}
  \textbf{Answer: (B).} Two-sided alternative requires doubling the tail probability.

  \item If the margin of error for a 1-sample $t$-interval is $\text{ME} = 3.5$, increasing the sample size will:
  \begin{enumerate}[label=(\Alph*)]
    \item Increase the margin of error \item Decrease the margin of error \item Have no effect on ME \item Increase the sample mean \item Make $t^*$ larger
  \end{enumerate}
  \textbf{Answer: (B).} Larger $n$ decreases standard error $s/\sqrt{n}$, reducing ME.
\end{enumerate}

\begin{enumerate}[resume, leftmargin=1.5em]
  \item Which of the following is an example of a matched pairs design?
  \begin{enumerate}[label=(\Alph*)]
    \item Comparing test scores of 30 girls and 30 boys.
    \item Comparing blood pressure measurements of 50 patients before and after taking a medication.
    \item Comparing corn yields from 10 plots in Texas and 10 plots in Iowa.
    \item Surveying 100 randomly chosen voters.
    \item Measuring weights of 40 dogs of different breeds.
  \end{enumerate}
  \textbf{Answer: (B).} Pre- and post-treatment measurements on the same individual constitute matched pairs.

  \item In a 1-sample $t$-test with $n = 20$, the test statistic is $t = -1.73$. For $H_a: \mu < \mu_0$, the $p$-value falls in which range?
  \begin{enumerate}[label=(\Alph*)]
    \item $p < 0.01$ \item $0.025 < p < 0.05$ \item $0.05 < p < 0.10$ \item $p > 0.10$ \item Exactly 0.05
  \end{enumerate}
  \textbf{Answer: (C).} For $df = 19$, $t_{0.05} = 1.729 \implies p \approx 0.050$. Specifically, $p$-value $\approx 0.0499 \approx 0.05$.

  \item If an investigator knows the true population standard deviation $\sigma = 15$, which procedure should be used to estimate $\mu$?
  \begin{enumerate}[label=(\Alph*)]
    \item 1-sample $t$-interval \item 1-sample $z$-interval \item 2-sample $t$-test \item Chi-square test \item Regression test
  \end{enumerate}
  \textbf{Answer: (B).} Knowing $\sigma$ allows using $z$ instead of $t$.

  \item What is the conservative degrees of freedom for a two-sample $t$-test with $n_1 = 15$ and $n_2 = 22$?
  \begin{enumerate}[label=(\Alph*)]
    \item 37 \item 35 \item 14 \item 21 \item 15
  \end{enumerate}
  \textbf{Answer: (C).} Conservative $df = \min(n_1 - 1, n_2 - 1) = \min(14, 21) = 14$.

  \item When constructing a confidence interval for $\mu$, the critical value $t^*$ depends on:
  \begin{enumerate}[label=(\Alph*)]
    \item Only the sample mean $\bar{x}$
    \item The confidence level and degrees of freedom ($n - 1$)
    \item The population size $N$
    \item The margin of error
    \item The $p$-value
  \end{enumerate}
  \textbf{Answer: (B).} $t^*$ is determined exclusively by the desired confidence level and $df$.

  \item If zero is contained within a 95\% confidence interval for $\mu_d$ in a paired experiment, what decision would be made for $H_0: \mu_d = 0$ at $\alpha = 0.05$?
  \begin{enumerate}[label=(\Alph*)]
    \item Reject $H_0$ \item Fail to reject $H_0$ \item Accept $H_a$ \item Increase sample size \item Conclude a Type I error was made
  \end{enumerate}
  \textbf{Answer: (B).} An interval containing 0 means the hypothesized value 0 is plausible, so we fail to reject $H_0$.

  \item In a sample of $n = 35$, a histogram shows slight right skewness. Can a $t$-test be conducted safely?
  \begin{enumerate}[label=(\Alph*)]
    \item No, data must be perfectly symmetric.
    \item Yes, because $n = 35 \ge 30$, the Central Limit Theorem ensures robustness.
    \item No, because $s$ cannot be calculated.
    \item Yes, only if $\mu = 0$.
    \item No, $z$ must be used instead.
  \end{enumerate}
  \textbf{Answer: (B).} $t$-procedures are robust to moderate skewness when $n \ge 30$.

  \item Standard error of the difference between two independent sample means is:
  \begin{enumerate}[label=(\Alph*)]
    \item $\sqrt{\frac{s_1^2}{n_1} + \frac{s_2^2}{n_2}}$ \item $\frac{s_1}{n_1} + \frac{s_2}{n_2}$ \item $\sqrt{s_1^2 + s_2^2}$ \item $\frac{s_1 + s_2}{\sqrt{n_1 + n_2}}$ \item $s_1 - s_2$
  \end{enumerate}
  \textbf{Answer: (A).} Standard formula for $\text{SE}(\bar{x}_1 - \bar{x}_2)$.

  \item If $n = 100, \bar{x} = 50, s = 10$, what is the margin of error for a 95\% confidence interval for $\mu$ ($t^* \approx 1.984$)?
  \begin{enumerate}[label=(\Alph*)]
    \item 1.984 \item 0.1984 \item 19.84 \item 1.00 \item 0.50
  \end{enumerate}
  \textbf{Answer: (A).} $\text{ME} = t^*(s/\sqrt{n}) = 1.984(10/\sqrt{100}) = 1.984(1.0) = 1.984$.

  \item Which of the following is true regarding $t$-distributions?
  \begin{enumerate}[label=(\Alph*)]
    \item They are skewed to the right.
    \item They have a larger variance than the standard normal distribution.
    \item They can take only positive values.
    \item They do not depend on sample size.
    \item They are used exclusively for categorical variables.
  \end{enumerate}
  \textbf{Answer: (B).} Fatter tails mean larger variance ($\text{Var} = df/(df-2) > 1$).
\end{enumerate}

\subsection*{Part II: Free-Response Practice Problem}

\begin{workbookbox}[FRQ 7.1: Paired $t$-Test for Matched Pairs]
\textbf{Problem:} A fitness coach tests whether a 6-week plyometric training program increases vertical jump height in basketball players. 12 players have their vertical jumps measured (in inches) before and after the program. The mean difference is $\bar{x}_d = \text{After} - \text{Before} = +2.4\text{ inches}$ with a sample standard deviation of differences $s_d = 1.8\text{ inches}$. A boxplot of the differences shows no outliers. Conduct a hypothesis test at $\alpha = 0.05$.

\textbf{Grader-Approved Score-4 Solution:}
\begin{itemize}[leftmargin=1.2em]
  \item \textbf{P:} Let $\mu_d = \text{the true mean difference (After } - \text{ Before) in vertical jump height for all players}$.
  \item \textbf{H:} $H_0: \mu_d = 0$ versus $H_a: \mu_d > 0$.
  \item \textbf{A:}
  1. \textit{Random:} Stated representative sample of players.
  2. \textit{10\% Condition:} $n = 12 \le 0.10(\text{all basketball players})$.
  3. \textit{Normal Condition:} $n = 12 < 30$, but the boxplot of differences displays no outliers or extreme skewness. Safe to proceed with $t$-distribution.
  \item \textbf{N:} Paired $t$-test for a mean difference $\mu_d$.
  \item \textbf{T:} $df = 12 - 1 = 11$.
  \[ t = \frac{\bar{x}_d - 0}{s_d / \sqrt{n}} = \frac{2.4 - 0}{1.8 / \sqrt{12}} = \frac{2.4}{0.5196} \approx 4.62 \]
  \item \textbf{O:} $p\text{-value} = P(t_{11} \ge 4.62) = \mathbf{0.00036}$.
  \item \textbf{M:} Because the $p$-value ($0.00036$) is strictly less than $\alpha = 0.05$, we reject $H_0$. There is convincing statistical evidence that the 6-week training program produces a true mean increase in vertical jump height.
\end{itemize}
\end{workbookbox}


\begin{frqrubricbox}[FRQ 7.2: Paired $t$-Test for Matched Pairs vs. Two-Sample $t$-Test]
\textbf{Problem:} An educational psychologist evaluates a spatial reasoning software module. Twelve high school seniors take a standardized 50-point spatial visualization test before and after completing a 4-week computer course.
The summary statistics of the test score differences ($d = \text{Post} - \text{Pre}$) are:
\[ n = 12, \quad \bar{x}_d = 6.50\text{ points}, \quad s_d = 2.40\text{ points} \]
A normal probability plot of the differences displays a strong linear pattern with no outliers.
\begin{enumerate}[label=(\alph*)]
  \item Explain why a paired $t$-test is appropriate for this data rather than a two-sample independent $t$-test.
  \item State the appropriate null and alternative hypotheses to test whether the software course significantly improves spatial visualization scores.
  \item Check all required conditions for conducting the paired $t$-test.
  \item Calculate the test statistic and $p$-value, and state your conclusion at the $\alpha = 0.01$ significance level.
\end{enumerate}

\textbf{Grader-Approved Score-4 Solution \& Rubric:}
\begin{itemize}[leftmargin=1.2em]
  \item \textbf{Part (a):} A paired $t$-test is appropriate because the data consist of matched pairs of measurements taken on the exact same subjects (pre-test and post-test scores for each student). The two sets of scores are inherently dependent, not two independent random samples.
  \item \textbf{Part (b):} Let $\mu_d$ be the true mean difference in spatial visualization scores ($\text{Post} - \text{Pre}$) for all students who complete the course.\\
  $H_0: \mu_d = 0 \quad \text{vs.} \quad H_a: \mu_d > 0$.
  \item \textbf{Part (c):}
  \textit{Random:} The 12 students are treated as representative of the target student population.\\
  \textit{10\% Condition:} $n = 12 \le 0.10(\text{All high school seniors})$.\\
  \textit{Normal/Large Sample:} $n = 12 < 30$, but the normal probability plot of differences is linear with no outliers, making the normality assumption plausible.
  \item \textbf{Part (d):}
  $\text{SE}(\bar{x}_d) = \frac{s_d}{\sqrt{n}} = \frac{2.40}{\sqrt{12}} \approx \frac{2.40}{3.4641} \approx 0.6928\text{ points}$.\\
  $t = \frac{\bar{x}_d - 0}{\text{SE}(\bar{x}_d)} = \frac{6.50}{0.6928} \approx 9.38$.\\
  Degrees of freedom: $df = n - 1 = 12 - 1 = 11$.\\
  $p$-value: $P(t_{11} \ge 9.38) < 0.00001$.\\
  \textit{Conclusion:} Because $p\text{-value} < 0.01$, we reject $H_0$. There is overwhelming statistical evidence that the spatial visualization software module significantly improves student test scores.
\end{itemize}
\end{frqrubricbox}

\begin{trapbox}[Unit 7 Top 5 AP Exam Pitfalls to Avoid]
\begin{enumerate}[leftmargin=1.2em]
  \item \textbf{Running a Two-Sample $t$-Test on Paired Data:} If subjects are matched or measured twice, you MUST run a paired $t$-test on differences ($df = n - 1$), NOT a two-sample test!
  \item \textbf{Selecting Pooled on the TI-84:} On the AP Statistics Exam, NEVER pool variances for a two-sample $t$-test. Always select \texttt{Pooled: NO}!
  \item \textbf{Degrees of Freedom for Two-Sample $t$:} Conservative $df = \min(n_1 - 1, n_2 - 1)$ is always accepted if TI-84 Satterthwaite $df$ is not used.
  \item \textbf{Checking Normality When $n < 30$:} When $n < 30$, you MUST check the plot of sample data for outliers and severe skewness.
  \item \textbf{Interpreting $p$-Value as Strength of Effect:} A tiny $p$-value means strong evidence against $H_0$; it does NOT mean the effect size is large!
\end{enumerate}
\end{trapbox}
\section{Unit 7 Master Summary}
\begin{lessonbox}[Unit 7 Essential Review Checklist]
\begin{itemize}[leftmargin=1.2em]
  \item \textbf{1-Sample $t$:} $\bar{x} \pm t^*(s/\sqrt{n})$, $df = n - 1$.
  \item \textbf{Normality Check for Means:} Normal pop OR $n \ge 30$ (CLT) OR graph with no strong skewness/outliers.
  \item \textbf{Paired $t$:} Run 1-sample $t$ on the list of differences ($d = x_1 - x_2$).
  \item \textbf{2-Sample $t$:} Independent groups. Never pool variances!
\end{itemize}
\end{lessonbox}
\clearpage
